% abstract.tex

\chapter*{Abstract}
\addcontentsline{toc}{chapter}{Abstract}

In Colombia, more than 314,000 families receive electricity through informal connections that meet no technical or legal standard. These settlements, known as subnormal neighborhoods (BS), face a problem that goes beyond infrastructure: structural poverty, distrust toward utility companies, and community resistance to billing have undermined normalization efforts for decades. This thesis proposes a comprehensive methodology for the normalization of electrical networks in BS that incorporates Small-Scale Self-Generation (AGPE) systems as part of the process, within the framework of the Electrical Network Normalization Program (PRONE), whose financing structure was modified in 2024 to include this type of investment.

The research begins with a detailed characterization of BS, with emphasis on the Caribbean region, which concentrates 90\% of the country's certified settlements. This draws on data from the Single Information System (SUI), the DANE Quality of Life Survey (2024), and the quinquennial investment plan reports of Network Operators (OR) Air-e and Afinia. The current regulatory framework for network normalization and AGPE systems was reviewed, and three case studies documented in Colombian literature were analyzed to identify what has worked, what has not, and why.

A distinctive element of this research is a direct consultation with 20 professionals from the electricity and public sectors with hands-on experience in normalization projects. Their responses revealed patterns that technical literature often misses: the main barrier to normalization is not technical but social — community resistance to individual metering and the perceived increase in electricity bills. This finding gives practical relevance to AGPE as a strategy for moderating that economic impact and facilitating community acceptance.

Based on a hypothetical case study built from actual conditions in the Caribbean region, the technical and financial viability of normalization projects with and without AGPE was assessed. Results show that the renewable generation component can significantly reduce the normalized user's monthly bill, making it both a technical and social argument for improving community buy-in.

The final output is a five-phase roadmap that guides OR from territory diagnosis through long-term asset sustainability, defining responsible actors, requirements, and deliverables for each stage. The 2024 procurement cycle, in which no project submitted to PRONE met eligibility requirements, confirms that this kind of methodological guidance is genuinely needed.

\bigskip
\noindent\textbf{Keywords:} subnormal neighborhoods, electrical network normalization, small-scale self-generation, PRONE, methodology, just energy transition, Caribbean region.